% GENERATED FILE. Edit canonical lessons, then run npm run generate:books.
% canonical-source-hash: fnv1a64:2b7f592f6eff2778
% canonical-lessons: ZH-C82-hou, ZH-C82-zuobian, ZH-C82-youbian, ZH-C82-zhongjian, ZH-C82-duimian

\chapter{Behind, Left and Right}
\label{ch:zh-behind-left-and-right}
\hlchaptermodality{82}

\begin{chapteropening}
\textbf{By the end of this chapter:} \emph{I can say after, left, right, in the middle and opposite.}

Complete the last lesson of chapter 82: I can say after, left, right, in the middle and opposite.
\end{chapteropening}

\section[hòu]{hòu — after; behind; later}

\label{lesson:ZH-C82-hou}



\begin{quote}
Before the new one: say the Mandarin for abroad, then the Mandarin for beside.
\end{quote}

\subsection*{You'll want to know: hòu}
\emph{hòu} — \textquotedblleft{}after; behind; later\textquotedblright{}.

\begin{quote}
\emph{sān tiān hòu wǒmen qù Běijīng} — In three days we're going to Beijing.
\end{quote}

Its characters come later; for now it is a word for the ear.

\begin{cousinweb}
\emph{hòu} is inherited from Middle Chinese \emph{huwX} (Guangyun, in Baxter's transcription), glossed with the pair 'before and after'; the dictionary writes it with the traditional form of the character.
\end{cousinweb}

\begin{grammarlens}[title={What you've built}]
One more word for meeting people.
\end{grammarlens}

\subsection*{Your turn}
\begin{itemize}
\raggedright
  \item \emph{Say it:} \emph{hòu}
  \item \emph{Say it:} \emph{hòu}, once more
  \item \emph{Say it:} \emph{sān tiān hòu wǒmen qù Běijīng} — In three days we're going to Beijing.
  \item \emph{Recall:} say the Mandarin for abroad, then the Mandarin for beside, then say \emph{hòu} again
\end{itemize}

\begin{tcolorbox}[breakable,colback=teal!4,colframe=teal!35!black,title={Before you move on}]
What does hòu mean? (\textquotedblleft{}After; behind; later\textquotedblright{}.) Say it once more.
\end{tcolorbox}
\section[zuǒbian]{zuǒbian — the left; on the left}

\label{lesson:ZH-C82-zuobian}



\begin{quote}
Before the new one: say the Mandarin for beside, then the Mandarin for after.
\end{quote}

\subsection*{You'll want to know: zuǒbian}
\emph{zuǒbian} — \textquotedblleft{}the left; on the left\textquotedblright{}.

\begin{quote}
\emph{zuǒbian shì wǒ de fángjiān} — My room is on the left.
\end{quote}

Its characters come later; for now it is a word for the ear.

\begin{cousinweb}
\emph{zuǒbian} is built from zuǒ 'left' and bian 'side', both inherited from Middle Chinese, \emph{tsaX} and \emph{pen} (Guangyun, in Baxter's transcription).
\end{cousinweb}

\begin{grammarlens}[title={What you've built}]
One more word for meeting people.
\end{grammarlens}

\subsection*{Your turn}
\begin{itemize}
\raggedright
  \item \emph{Say it:} \emph{zuǒbian}
  \item \emph{Say it:} \emph{zuǒbian}, once more
  \item \emph{Say it:} \emph{zuǒbian shì wǒ de fángjiān} — My room is on the left.
  \item \emph{Recall:} say the Mandarin for beside, then the Mandarin for after, then say \emph{zuǒbian} again
\end{itemize}

\begin{tcolorbox}[breakable,colback=teal!4,colframe=teal!35!black,title={Before you move on}]
What does zuǒbian mean? (\textquotedblleft{}The left; on the left\textquotedblright{}.) Say it once more.
\end{tcolorbox}
\section[yòubian]{yòubian — the right; on the right}

\label{lesson:ZH-C82-youbian}



\begin{quote}
Before the new one: say the Mandarin for after, then the Mandarin for the left side.
\end{quote}

\subsection*{You'll want to know: yòubian}
\emph{yòubian} — \textquotedblleft{}the right; on the right\textquotedblright{}.

\begin{quote}
\emph{yòubian shì shāngdiàn} — The shop is on the right.
\end{quote}

Its characters come later; for now it is a word for the ear.

\begin{cousinweb}
\emph{yòubian} is built from yòu 'right' and bian 'side', both inherited from Middle Chinese, \emph{hjuwX} and \emph{pen} (Guangyun, in Baxter's transcription).
\end{cousinweb}

\begin{grammarlens}[title={What you've built}]
One more word for meeting people.
\end{grammarlens}

\subsection*{Your turn}
\begin{itemize}
\raggedright
  \item \emph{Say it:} \emph{yòubian}
  \item \emph{Say it:} \emph{yòubian}, once more
  \item \emph{Say it:} \emph{yòubian shì shāngdiàn} — The shop is on the right.
  \item \emph{Recall:} say the Mandarin for after, then the Mandarin for the left side, then say \emph{yòubian} again
\end{itemize}

\begin{tcolorbox}[breakable,colback=teal!4,colframe=teal!35!black,title={Before you move on}]
What does yòubian mean? (\textquotedblleft{}The right; on the right\textquotedblright{}.) Say it once more.
\end{tcolorbox}
\section[zhōngjiān]{zhōngjiān — the middle; between}

\label{lesson:ZH-C82-zhongjian}



\begin{quote}
Before the new one: say the Mandarin for the left side, then the Mandarin for the right side.
\end{quote}

\subsection*{You'll want to know: zhōngjiān}
\emph{zhōngjiān} — \textquotedblleft{}the middle; between\textquotedblright{}.

\begin{quote}
\emph{zhōngjiān nàge shì wǒ} — The one in the middle is me.
\end{quote}

Its characters come later; for now it is a word for the ear.

\begin{cousinweb}
\emph{zhōngjiān} is built from \zh{中} 'middle' and jiān 'a gap', both inherited from Middle Chinese, \emph{trjuwng} and \emph{kean} (Guangyun, in Baxter's transcription); the dictionary lists the second under an older form of the character.
\end{cousinweb}

\begin{grammarlens}[title={What you've built}]
One more word for meeting people.
\end{grammarlens}

\subsection*{Your turn}
\begin{itemize}
\raggedright
  \item \emph{Say it:} \emph{zhōngjiān}
  \item \emph{Say it:} \emph{zhōngjiān}, once more
  \item \emph{Say it:} \emph{zhōngjiān nàge shì wǒ} — The one in the middle is me.
  \item \emph{Recall:} say the Mandarin for the left side, then the Mandarin for the right side, then say \emph{zhōngjiān} again
\end{itemize}

\begin{tcolorbox}[breakable,colback=teal!4,colframe=teal!35!black,title={Before you move on}]
What does zhōngjiān mean? (\textquotedblleft{}The middle; between\textquotedblright{}.) Say it once more.
\end{tcolorbox}
\section[duìmiàn]{duìmiàn — opposite; across from}

\label{lesson:ZH-C82-duimian}



\begin{quote}
Before the new one: say the Mandarin for the right side, then the Mandarin for the middle.
\end{quote}

\subsection*{You'll want to know: duìmiàn}
\emph{duìmiàn} — \textquotedblleft{}opposite; across from\textquotedblright{}.

\begin{quote}
\emph{wǒ jiā duìmiàn yǒu yí ge fàndiàn} — There's a restaurant across from my home.
\end{quote}

Its characters come later; for now it is a word for the ear.

\begin{cousinweb}
\emph{duìmiàn} is built from \zh{对} 'to face' and miàn 'face', both inherited from Middle Chinese, \emph{twojH} and \emph{mjienH} (Guangyun, in Baxter's transcription): face to face.
\end{cousinweb}

\begin{grammarlens}[title={What you've built}]
One more word for meeting people.
\end{grammarlens}

\subsection*{Your turn}
\begin{itemize}
\raggedright
  \item \emph{Say it:} \emph{duìmiàn}
  \item \emph{Say it:} \emph{duìmiàn}, once more
  \item \emph{Say it:} \emph{wǒ jiā duìmiàn yǒu yí ge fàndiàn} — There's a restaurant across from my home.
  \item \emph{Recall:} say the Mandarin for the right side, then the Mandarin for the middle, then say \emph{duìmiàn} again
\end{itemize}

\begin{tcolorbox}[breakable,colback=teal!4,colframe=teal!35!black,title={Before you move on}]
What does duìmiàn mean? (\textquotedblleft{}Opposite; across from\textquotedblright{}.) Say it once more.
\end{tcolorbox}
